% TEMPLATE-MILC-v3.tex - plantilla maestra UNICA de las guias MILC v3.
%
% La usan las cuatro familias de guias, cada una con su driver:
%   · Grados 6-11          -> scripts/build-guias-g11.py
%   · Bebras               -> scripts/build-guias-bebras.py
%   · Semillero            -> scripts/build-guias-semillero.py
%   · Territorio interior  -> scripts/build-guias-territorio-interior.py
%
% Antes habia tres copias casi identicas (~400 lineas iguales, 6 distintas):
% cada mejora habia que aplicarla tres veces, y en la practica no se hacia
% (el motor de recursos vivia solo en dos de ellas). Lo que varia entre
% familias esta parametrizado en 6 placeholders que cada driver llena
% (se nombran aqui SIN su sintaxis real: el builder reemplaza por texto plano,
% tambien dentro de los comentarios, y un valor multilinea rompe el preambulo):
%
%   HEADER_IZQ        encabezado izquierdo de pagina
%   HEADER_DER        encabezado derecho de pagina
%   PORTADA_ETIQUETA  rotulo sobre el numero grande (GRADO/MOMENTO/MODULO)
%   PORTADA_SERIE     linea de serie de la portada
%   PORTADA_ID        identificador de la guia en la portada
%   PORTADA_CIRCULO   el circulito de grado (°), o vacio si no aplica
%   PDF_TITULO        titulo en texto plano (sin \textbf ni comillas TeX) para
%                     los metadatos del PDF (/Title); lo produce lib_xelatex.texto_pdf
%
% Composicion (auditoria 2026-09-03): las bandas de estacion piden espacio
% (\Needspace*) y prohiben el salto justo despues (\nopagebreak), asi no quedan
% huerfanas al pie; las cajas largas (softbox, drawbox, infoband, puentebox) son
% `breakable`, asi una caja de media pagina ya no arrastra todo a la siguiente
% dejando la anterior al 40 %. Todo texto sobre mostaza o verde va en milcNegro
% (blanco sobre #E5B400 da 1,93:1 y sobre #58A923 2,95:1; ninguno pasa WCAG AA).
% El resto de claves las documenta content/guias/_SCHEMA.md. El script valida
% que no queden placeholders sin reemplazar antes de escribir el archivo final.

% Etiquetado PDF (latex-lab): /Lang, estructura y texto alternativo de las
% figuras, para lectores de pantalla. Debe ir ANTES de \documentclass.
\DocumentMetadata{lang=es-CO,tagging=on}
\documentclass[11pt,letterpaper]{article}
% headheight=24pt: la cabecera derecha lleva dos lineas (identidad de la guia
% y estacion actual). headsep se acorta en la misma medida para que el bloque
% de texto quede exactamente donde estaba con headheight=12pt.
\usepackage[margin=2.0cm,headheight=24pt,headsep=13pt]{geometry}
\usepackage[table]{xcolor}
\usepackage{fontspec}
\usepackage[spanish]{babel}
\usepackage{tikz}
% Librerías TikZ para los diagramas de las guías (bloque `recursos:`):
%  · babel  → babel en español vuelve activos caracteres como «>», y TikZ los usa
%             en sus opciones (>=stealth, ->). Sin esta librería, cualquier
%             diagrama con flechas revienta con «\language@active@arg> has an
%             extra }». Debe ir ANTES de que se dibuje nada.
%  · positioning        → sintaxis «below left=2pt and 3pt».
%  · decorations.pathreplacing → llaves ({brace}) para señalar rangos.
\usetikzlibrary{babel,positioning,decorations.pathreplacing}
\usepackage{graphicx}
% Circuitos: el trabajo de electrónica se hace hoy con micro:bit y sus sensores
% integrados (MakeCode, sin cableado), así que no se necesitan esquemáticos.
% Si algún día entran montajes con protoboard, basta descomentar esta línea:
% los diagramas .tex/.tikz declarados en `recursos:` ya se incrustan solos.
% \usepackage{circuitikz}
\usepackage[most]{tcolorbox}
\usepackage{tabularx}
\usepackage{array}
\usepackage{fancyhdr}
\usepackage{titlesec}
\usepackage[document]{ragged2e}  % bandera a la derecha en todo el cuerpo
\usepackage{enumitem}
\usepackage{microtype}
\usepackage{etoolbox}
\usepackage{needspace}
% hyperref al final del preambulo: metadatos del PDF (titulo, autor, idioma,
% familia), marcadores por estacion (\pdfbookmark) y enlaces sin recuadro.
\usepackage[hidelinks,
  pdftitle={Derechos de autor y propiedad intelectual con IA — Ley 23 de 1982 y obras derivadas},
  pdfauthor={PhD. Álvaro Cárdenas Orozco},
  pdfsubject={Tecnología e Informática · Grado 10 · Guía 6 · MILC v3},
  pdflang={es-CO},
  bookmarksopen=true,bookmarksnumbered=false]{hyperref}

% Helvetica Neue no trae la flecha U+2192, asi que cada «→» escrito literalmente
% en el YAML se compilaba a NADA, en silencio (462 flechas en 61 guias: rutas del
% tipo «paleta → tipografia → lockup» salian rotas). Mapeamos el caracter a su
% version matematica, que si tiene glifo. Asi el autor puede seguir escribiendo
% «→» comodamente en el YAML.
\usepackage{newunicodechar}
\newunicodechar{→}{\ensuremath{\rightarrow}}
% Mismo problema con los simbolos de marca y vinetas: el «✓» de «una palomita ✓»
% salia como cuadrito vacio en el PDF de todas las guias que lo usaban.
\usepackage{pifont}
\newunicodechar{✓}{\ding{51}}
\newunicodechar{✗}{\ding{55}}
\newunicodechar{✔}{\ding{52}}
\newunicodechar{•}{\textbullet}
\newunicodechar{≥}{\ensuremath{\geq}}
\newunicodechar{≤}{\ensuremath{\leq}}

\setmainfont{Helvetica Neue}
\setsansfont{Helvetica Neue}
\setlength{\parindent}{0pt}
\setlength{\parskip}{6pt}
% Sin particion de palabras: en 6.º-7.º una palabra cortada por guion al final
% de linea («Comuni-cacion») frena la lectura mucho mas que una linea corta.
\hyphenpenalty=10000
\exhyphenpenalty=10000
\pretolerance=2000
\tolerance=3000
\setlength{\emergencystretch}{3em}
\linespread{1.10}
\renewcommand{\arraystretch}{1.25}
\setlist{leftmargin=5mm,itemsep=1mm,topsep=1mm}
\newcolumntype{Y}{>{\RaggedRight\arraybackslash}X}

\definecolor{milcMagenta}{HTML}{E6007E}
\definecolor{milcVino}{HTML}{5A0038}
\definecolor{milcTurquesa}{HTML}{0093A5}
\definecolor{milcVerde}{HTML}{58A923}
\definecolor{milcMostaza}{HTML}{E5B400}
\definecolor{milcOcre}{HTML}{B86600}
\definecolor{milcNegro}{HTML}{241816}
\definecolor{milcGris}{HTML}{F7F7F4}
\definecolor{milcLinea}{HTML}{E8E2DD}
\definecolor{milcGradePrimary}{HTML}{FF2D87}
\definecolor{milcGradeDark}{HTML}{9F1B56}
\definecolor{milcGradeSoft}{HTML}{FFE0EE}
\definecolor{milcGradeAccent}{HTML}{CFE7FF}
\definecolor{milcGradeText}{HTML}{FFFFFF}
\color{milcNegro}

\pagestyle{fancy}
\fancyhf{}
% Cabecera de dos lineas: identidad de la guia arriba y, a la derecha y en
% gris, la estacion en curso (\markright desde \milcestacion). Antes de la
% primera estacion la segunda linea va vacia; el \strut mantiene la altura
% para que la linea de arriba no baile entre paginas.
\lhead{\small Tecnología e Informática · Grado 10\\[-1pt]{\footnotesize\strut}}
\rhead{\small Guía 6 · MILC v3\\[-1pt]{\footnotesize\strut\color{milcNegro!65}\rightmark}}
\cfoot{\small\thepage}
\renewcommand{\headrulewidth}{0pt}

% Sobre mostaza (#E5B400) y verde (#58A923) el texto va en milcNegro; sobre el
% resto de la paleta, en blanco. Se usa dentro de fonttitle/fontupper, donde un
% \color posterior manda sobre coltitle.
\newcommand{\milctintasobre}[1]{%
  \ifboolexpr{test{\ifstrequal{#1}{milcMostaza}} or test{\ifstrequal{#1}{milcVerde}}}%
    {\color{milcNegro}}{\color{white}}}

\newtcolorbox{titlebox}[1]{
  enhanced,colback=#1,colframe=#1,arc=7mm,boxrule=0pt,
  left=7mm,right=7mm,top=3mm,bottom=3mm,
  before skip=8pt,after skip=8pt,
  fontupper=\sffamily\bfseries\Large\color{white}
}

\newtcolorbox{softbox}[3]{
  enhanced,breakable,pad at break=3mm,
  colback=#2,colframe=#1,borderline west={4pt}{0pt}{#1},
  arc=2mm,boxrule=.4pt,left=4mm,right=4mm,top=3mm,bottom=3mm,
  before skip=6pt,after skip=8pt,title={#3},coltitle=white,
  colbacktitle=#1,fonttitle=\sffamily\bfseries\milctintasobre{#1},fontupper=\RaggedRight,
  attach boxed title to top left={xshift=3mm,yshift=-2mm},
  boxed title style={arc=2mm, boxrule=0pt}
}

\newtcolorbox{drawbox}[2]{
  enhanced,breakable,pad at break=4mm,
  colback=white,colframe=#1,arc=4mm,boxrule=1pt,
  left=5mm,right=5mm,top=4mm,bottom=4mm,before skip=8pt,after skip=8pt,
  title={#2},coltitle=white,colbacktitle=#1,fonttitle=\sffamily\bfseries\milctintasobre{#1},
  fontupper=\RaggedRight,
  attach boxed title to top left={xshift=4mm,yshift=-2mm},
  boxed title style={arc=3mm, boxrule=0pt}
}

\newtcolorbox{infoband}[1]{
  enhanced,breakable,pad at break=2mm,colback=#1!8,colframe=#1!50,arc=2mm,boxrule=.5pt,
  left=4mm,right=4mm,top=2mm,bottom=2mm,before skip=4pt,after skip=4pt,
  fontupper=\small\RaggedRight
}

% Puente narrativo entre secciones (contrato editorial Conéctate).
% Da continuidad: "Acabas de X. Ahora vas a Y porque Z."
% Visualmente: banda izquierda en color del grado, texto en itálica pequeña.
\newtcolorbox{puentebox}{
  enhanced,breakable,pad at break=2mm,colback=milcGradeSoft,colframe=milcGradeDark,
  borderline west={3pt}{0pt}{milcGradeDark},
  arc=0mm,boxrule=0pt,left=5mm,right=4mm,top=2.5mm,bottom=2.5mm,
  before skip=4pt,after skip=8pt,
  fontupper=\itshape\small\color{milcNegro}\RaggedRight
}
% La flecha va en modo matematico a proposito: Helvetica Neue NO tiene el glifo
% U+2192, asi que \textrightarrow se compilaba a NADA («Missing character: There
% is no → in font Helvetica Neue») y el puente salia sin flecha. En modo
% matematico la flecha viene de la fuente matematica y si se dibuja.
\newcommand{\puente}[1]{%
  \begin{puentebox}%
    \textcolor{milcGradeDark}{$\rightarrow$}\hspace{2mm}#1%
  \end{puentebox}%
}

\newcommand{\linead}[1]{\noindent\color{milcLinea}\rule{#1}{0.5pt}\color{milcNegro}}
\newcommand{\respuesta}[1]{\par\vspace{2mm}\foreach \n in {1,...,#1}{\noindent\color{milcLinea}\rule{\linewidth}{0.5pt}\par\vspace{5mm}}\color{milcNegro}}
\newcommand{\checkbox}{\textcolor{milcGradeDark}{\fbox{\phantom{x}}}\hspace{5pt}}

% ─── Listas aireadas para las guías socioemocionales ────────────────────
% Componer las verificaciones como «\checkbox texto\\» deja el texto que salta
% de línea debajo del cuadrito y apelmaza el bloque. Estos entornos alinean la
% continuación y airean los ítems. Los usa Territorio Interior; cualquier
% familia puede hacerlo.
\newlist{checklista}{itemize}{1}
\setlist[checklista]{
  label=\textcolor{milcGradeDark}{\fbox{\phantom{x}}},
  leftmargin=9mm, labelsep=3.2mm, itemsep=2.4mm,
  topsep=2.5mm, parsep=0pt, partopsep=0pt, before=\RaggedRight
}
\newlist{pasoslista}{enumerate}{1}
\setlist[pasoslista]{
  label=\textcolor{milcGradeDark}{\sffamily\bfseries\arabic*.},
  leftmargin=8mm, labelsep=3mm, itemsep=2.8mm,
  topsep=1mm, parsep=0pt, partopsep=0pt, before=\RaggedRight
}

\newcommand{\phasecard}[4]{
\begin{tcolorbox}[enhanced,colback=#3,colframe=#2,arc=3mm,boxrule=.5pt,height=3.05cm,valign=center,halign=center,left=1.5mm,right=1.5mm,top=2mm,bottom=2mm]
{\sffamily\bfseries\fontsize{22}{24}\selectfont\textcolor{#2}{#1}}\par
{\sffamily\bfseries\small #4}
\end{tcolorbox}
}

% ── Piezas del rediseno 2026 (legibilidad 6.º-11.º) ───────────────────
% Banda de estacion: reemplaza el titlebox macizo. Titulo a la izquierda,
% dato practico (tiempo/modalidad) a la derecha, en una sola linea.
% Nunca huerfana: pide 9 lineas libres (o salta de pagina rellenando con
% \vfill) y prohibe el salto justo despues de la banda. Nueve y no seis:
% la caja partible que sigue necesita ~80pt (titulo adosado, relleno y dos
% lineas) para arrancar en la misma pagina; con menos, tcolorbox la manda
% entera a la siguiente y la banda se queda sola. El penalty 10000 va
% en `after`, ANTES del espacio posterior: un `after skip` (aunque sea 0pt)
% es un \vskip tras la caja, y ese glue es un punto de corte legal que un
% \nopagebreak escrito despues de \end{tcolorbox} ya no cierra. Cada banda
% deja marcador PDF y actualiza la cabecera derecha.
\newcounter{milcestacion}
\newcommand{\milcestacion}[2]{%
\Needspace*{9\baselineskip}%
\stepcounter{milcestacion}%
\pdfbookmark[1]{#2}{milcestacion\themilcestacion}%
\markright{#2}%
\begin{tcolorbox}[enhanced,colback=#1,colframe=#1,arc=2mm,boxrule=0pt,
  left=5mm,right=5mm,top=2.6mm,bottom=2.6mm,before skip=11pt,
  after={\par\nopagebreak\vskip7pt}]
{\sffamily\bfseries\fontsize{15}{18}\selectfont\milctintasobre{#1}#2}
\end{tcolorbox}}

% Tarjeta de paso numerada: sustituye las tablas «Pilar / Aplicacion», que
% eran formato de consulta docente, no de estudiante. El numero va en un
% circulo sobre el borde para que la secuencia se lea de un vistazo.
\newcommand{\milcpaso}[3]{%
\Needspace{4\baselineskip}%
\begin{tcolorbox}[enhanced,colback=white,colframe=milcLinea,arc=1.5mm,boxrule=.9pt,
  left=14mm,right=4mm,top=3mm,bottom=3mm,before skip=4pt,after skip=4pt,
  fontupper=\RaggedRight,
  overlay={\node[anchor=north west,circle,fill=#2,text=white,inner sep=0pt,
    minimum size=7.4mm,font=\sffamily\bfseries\small]
    at ([xshift=3.6mm,yshift=-3.2mm]frame.north west) {#1};}]
#3
\end{tcolorbox}}

% Casilla marcable de tamano util con lapiz (la anterior era \fbox{\phantom{x}},
% de alto variable segun el texto que la seguia).
\newcommand{\milcmarca}{\framebox[3.8mm]{\rule{0pt}{3.4mm}}}

% Fila marcable de lista suelta (checks de la estacion 1).
\newcommand{\milccheckrow}[1]{%
\par\vspace{1.8mm}\noindent
\makebox[6.4mm][l]{\raisebox{-.55mm}{\textcolor{milcNegro}{\milcmarca}}}%
\parbox[t]{\dimexpr\linewidth-6.4mm\relax}{\RaggedRight #1}\par}

% Celda marcable dentro de la rubrica (tabla de autoevaluacion). Misma
% sangria francesa, pero pensada para vivir dentro de una columna X.
\newcommand{\milcopcion}[1]{%
\makebox[5.2mm][l]{\raisebox{-.55mm}{\textcolor{milcNegro}{\milcmarca}}}%
\parbox[t]{\dimexpr\linewidth-5.2mm\relax}{\RaggedRight #1}}

% ── Recursos visuales de la guía ──────────────────────────────────────
% Declarados en el bloque `recursos:` del YAML y emitidos por el builder.
% \guiaFigura   → imagen rasterizada o PDF (foto, carta celeste, montaje).
% \guiaDiagrama → fuente TikZ/circuitikz (.tex) incrustada como vector:
%                 es la vía para circuitos y esquemáticos editables.
% [#1] = texto alternativo (alt) · #2 = ruta relativa al .tex · #3 = pie
% El alt llega al PDF etiquetado (/Alt) y lo lee el lector de pantalla.
\newcommand{\guiaFigura}[3][]{%
\begin{center}
\includegraphics[width=0.88\linewidth,height=0.38\textheight,keepaspectratio,alt={#1}]{#2}\\[1.5mm]
{\footnotesize\itshape\textcolor{milcNegro!75}{#3}}
\end{center}
}

\newcommand{\guiaDiagrama}[2]{%
\begin{center}
\input{#1}\\[1.5mm]
{\footnotesize\itshape\textcolor{milcNegro!75}{#2}}
\end{center}
}

\begin{document}
\thispagestyle{empty}

% ── Portada ───────────────────────────────────────────────────────────
% Antes: un rectangulo de color a pagina completa. Se veia bien en pantalla,
% pero en la fotocopiadora del colegio se comia el toner y salia gris sucio.
% Ahora la banda ocupa el tercio superior y el resto va en blanco.
\begin{tikzpicture}[remember picture,overlay]
  \path[fill=milcGradePrimary]
    (current page.north west) rectangle ([yshift=-10.6cm]current page.north east);

  \node[anchor=north west,text=milcGradeText] at ([xshift=2.0cm,yshift=-1.55cm]current page.north west)
    {\sffamily\bfseries\fontsize{12}{14}\selectfont GRADO};
  \node[anchor=north west,text=milcGradeText,opacity=.85] at ([xshift=2.0cm,yshift=-2.35cm]current page.north west)
    {\sffamily\bfseries\fontsize{8.5}{10}\selectfont SERIE GUÍAS · TECNOLOGÍA E INFORMÁTICA};
  \node[anchor=north east,text=milcGradeText,opacity=.85,align=right] at ([xshift=-2.0cm,yshift=-1.55cm]current page.north east)
    {\sffamily\bfseries\fontsize{9}{12}\selectfont I.E. Sor María Juliana\\Cartago, Valle del Cauca};

  \node[anchor=west,text=milcGradeText] (gradonum) at ([xshift=1.85cm,yshift=-7.1cm]current page.north west)
    {\sffamily\bfseries\fontsize{112}{112}\selectfont 10};
  % El circulito de grado (°) solo aplica a las guias de grado («11°»); en
  % Bebras y Semillero el numero grande es un momento/modulo, no un grado.
  % Se posiciona relativo al nodo `gradonum`, no en coordenadas absolutas.
  \draw[milcGradeText,line width=1.6pt]
    ([xshift=2.0mm,yshift=-4.5mm]gradonum.north east) circle (.15cm);

  \node[anchor=west,text=milcGradeText,text width=11.4cm,align=left] at ([xshift=1.5cm,yshift=0.5cm]gradonum.east)
    {
     {\sffamily\bfseries\fontsize{9}{11}\selectfont\MakeUppercase{Período 1 · Escritura abierta con IA}}\\[3mm]
     {\sffamily\bfseries\fontsize{21}{25}\selectfont\setlength{\baselineskip}{27pt}Derechos de autor\\[4pt]con IA --- Ley 23\\[4pt]de 1982}\\[3.5mm]
     {\sffamily\bfseries\fontsize{15}{18}\selectfont Guía 6-10-TIC}
    };
\end{tikzpicture}

\vspace*{9.4cm}
\vfill

{\sffamily\bfseries\fontsize{8}{10}\selectfont\textcolor{milcGradeDark}{LO QUE TE LLEVAS HOY}}

\begin{tcolorbox}[enhanced,colback=white,colframe=milcNegro,arc=2mm,boxrule=1.6pt,
  left=5mm,right=5mm,top=4mm,bottom=4mm,before skip=3pt,after skip=12pt,
  fontupper=\RaggedRight\fontsize{11}{15.5}\selectfont]
Ficha de derechos de tu libro con autoría, licencia justificada, declaración explícita del uso de IA y fuentes atribuidas, más el análisis de 3 casos recientes
\end{tcolorbox}

\begin{tcolorbox}[enhanced,colback=milcGris,colframe=milcGris,arc=2mm,boxrule=0pt,
  left=5mm,right=5mm,top=3.5mm,bottom=3.5mm,after skip=12pt]
{\sffamily\bfseries\fontsize{8}{10}\selectfont\textcolor{milcNegro!55}{ESTA GUÍA ES DE}}\\[3mm]
\begin{tabularx}{\linewidth}{X p{3.2cm} p{3.2cm}}
\linead{\linewidth} & \linead{3.0cm} & \linead{3.0cm}\\[-1mm]
{\fontsize{8.5}{10}\selectfont\textcolor{milcNegro!55}{Nombre completo}} &
{\fontsize{8.5}{10}\selectfont\textcolor{milcNegro!55}{Grupo}} &
{\fontsize{8.5}{10}\selectfont\textcolor{milcNegro!55}{Fecha}}\\
\end{tabularx}
\end{tcolorbox}

\vfill

\noindent\textcolor{milcLinea}{\rule{\linewidth}{1pt}}\\[2mm]
{\sffamily\bfseries\fontsize{9.5}{12}\selectfont Autor: Dr. Álvaro Cárdenas Orozco}\\[1mm]
{\sffamily\fontsize{8.5}{11}\selectfont\textcolor{milcNegro!55}{Metodología MILC · apertura ancestral · pensamiento computacional · triángulo Dussel–estoicismo–Floridi}}

\newpage

\section*{Apertura: Derechos de autor y propiedad intelectual con IA --- Ley 23 de 1982 y obras derivadas}

\begin{softbox}{milcGradeDark}{milcGradeSoft}{Saber ancestral}
En 1891 el Estado colombiano compró 122 piezas de oro hechas por orfebres prehispánicos del Cauca medio. Dos años después, el 4 de mayo de 1893, las regaló a la reina regente de España. En 2017 la Corte Constitucional ordenó repatriarlas o readquirirlas, por vulneración de la moralidad administrativa y de la defensa del patrimonio público (Corte Constitucional de Colombia, 2017). Las piezas siguen en Madrid. La \textbf{cara de exclusión} es doble. «Quimbaya» es un exónimo: el nombre se lo pusieron los coleccionistas, no quienes hicieron las piezas. Y esos orfebres no tienen voz en el pleito. Lo que se discute es entre Estados, sobre quién es dueño de lo que hicieron otros.\par\smallskip{\footnotesize\textcolor{milcNegro!70}{\textbf{Fuente.} Corte Constitucional de Colombia. (2017). \emph{Sentencia SU-649/17} (M. P. Alberto Rojas Ríos), 19 de octubre.}}
\end{softbox}

\begin{softbox}{milcTurquesa}{milcGris}{Saber contemporáneo}
En Colombia los \textbf{derechos de autor} se rigen por la Ley 23 de 1982 y la Decisión Andina 351 de 1993. Hay dos tipos. Los \textbf{derechos morales} son irrenunciables, intransferibles y perpetuos: ser reconocido como autor, que nadie altere la obra, decidir si se publica. Nunca dejan de ser tuyos. Los \textbf{derechos patrimoniales} sí se pueden ceder: reproducir, distribuir, traducir, adaptar, cobrar. Sobre IA hay tres consensos. Un modelo no es autor, porque no es una persona. Si la intervención humana es sustancial, la obra es de quien la hizo. Y declarar el uso de la herramienta es práctica profesional, aunque la ley todavía no lo exija.
\end{softbox}

\begin{softbox}{milcMostaza}{milcGris}{Pregunta-puente}
\textit{¿De quién es una pieza cuando quienes la hicieron no tienen voz en el pleito? ¿Y de quién es un capítulo que una herramienta armó a partir de textos de miles de autores?}
\end{softbox}

\begin{softbox}{milcVerde}{milcGris}{Saber y saber hacer}
Al terminar podrás: (1) \textbf{identificar} en un libro real quién firma, bajo qué licencia circula y qué declara sobre su proceso; (2) \textbf{analizar} la diferencia entre derechos morales y patrimoniales, y tres casos recientes de disputa entre IA y derechos; (3) \textbf{crear} la ficha de derechos de tu libro, con la licencia justificada y el uso de IA declarado.
\end{softbox}



\small
\begin{tabularx}{\textwidth}{>{\bfseries}p{3.0cm}Y}
\rowcolor{milcGradePrimary!12}Autor & Dr. Álvaro Cárdenas Orozco\\
DBA & Producción editorial mediada por IA con criterio ético y técnico (MEN, T\&I)\\
Referentes & Ley 23 de 1982 · Floridi (ética de la IA generativa) · Dussel (autoría con dignidad)\\
Producto final & Ficha de derechos de tu libro con autoría, licencia justificada, declaración explícita del uso de IA y fuentes atribuidas, más el análisis de 3 casos recientes\\
\end{tabularx}
\normalsize

\puente{En la sesión 5 dejaste el plano del libro. Hoy defines quién lo firma y bajo qué condiciones circula. Es la sesión donde el libro deja de ser un ejercicio y se vuelve una obra con dueño declarado.}

\milcestacion{milcGradeDark}{Ruta MILC: así vamos a aprender}
\begin{tabularx}{\textwidth}{>{\centering\arraybackslash}X>{\centering\arraybackslash}X>{\centering\arraybackslash}X>{\centering\arraybackslash}X}
\phasecard{1}{milcVerde}{milcGris}{Escucha\\[-1mm]\small Observo, escucho y pregunto.} &
\phasecard{2}{milcTurquesa}{milcGris}{Entender\\[-1mm]\small Sistematizo y ordeno.} &
\phasecard{3}{milcMagenta}{milcGris}{Hacer\\[-1mm]\small Construyo y aplico.} &
\phasecard{4}{milcVino}{milcGris}{Cierre\\[-1mm]\small Evalúo y reflexiono.}
\end{tabularx}


\puente{Antes de teorizar, vas a abrir un libro real por la página de créditos. Ahí está, en letra pequeña, todo lo que hoy vas a aprender a escribir para el tuyo.}

\milcestacion{milcVerde}{Estación 1 de 3 · Escucha}

\begin{softbox}{milcVerde}{milcGris}{Escena para escuchar}
\textbf{Actividad 1 · IDENTIFICA --- Un libro y sus créditos} (15 min · individual). (1) Toma un libro que tengas a mano, en papel o digital. (2) Busca la página de créditos, que suele estar después de la portadilla. (3) Anota quién es el autor declarado con nombre completo. (4) Anota la editorial, el año y el país de edición. (5) Anota qué dice sobre derechos: si están todos reservados, si hay una licencia abierta, y si menciona traducción, ilustración o cualquier otra mano.
\end{softbox}

{\sffamily\bfseries\fontsize{8}{10}\selectfont\textcolor{milcNegro!55}{MARCA CUANDO LO HAYAS HECHO}}
\milccheckrow{Revisé la página de créditos de un libro real, no la portada.}
\milccheckrow{Anoté autoría, editorial, año y qué dice sobre derechos.}
\milccheckrow{Anoté si aparecen otras manos: traducción, ilustración, edición.}

\begin{infoband}{milcVerde}


\textbf{Lo que va al cuaderno (Actividad 1 · IDENTIFICA).} \textbf{Título:} «Actividad 1 --- Un libro y sus créditos». \textbf{Formato:} ficha con autoría, editorial, año, qué dice sobre derechos y qué otras manos aparecen nombradas. \textbf{Extensión:} un tercio de página. \textbf{Sabes que terminaste cuando} reconoces que cada libro trae una firma legal explícita, aunque casi nadie la lea.
\end{infoband}

\puente{Ya viste una página de créditos. Ahora vas a entender la \textbf{anatomía de los derechos}: qué es intransferible y qué se puede ceder, qué ofrece cada licencia y qué se sabe hoy sobre IA y autoría.}

\milcestacion{milcTurquesa}{Estación 2 de 3 · Entender}

\begin{softbox}{milcTurquesa}{milcGris}{Lo que vas a entender}
\textbf{Actividad 2 · ANALIZA --- Morales, patrimoniales y licencias} (20 min · parejas). (1) Con tu pareja, escriban los dos tipos de derecho con un ejemplo de cada uno. (2) Escriban las licencias abiertas más usadas y qué permite cada una. (3) Escriban los tres consensos sobre IA y autoría. (4) Cada uno decide qué licencia le conviene a su libro y le explica al otro por qué, y el otro le busca una objeción.
\end{softbox}

\milcpaso{1}{milcTurquesa}{Los dos tipos de derecho se descomponen así. Los \textbf{morales} son irrenunciables, intransferibles y perpetuos: la paternidad de la obra, su integridad y la decisión de publicarla. Siguen siendo tuyos aunque cedas todo lo demás. Los \textbf{patrimoniales} sí se ceden: reproducción, distribución, traducción, adaptación y cobro. Son los que se negocian con una editorial.}
\milcpaso{2}{milcTurquesa}{Las licencias abiertas más usadas parten todas de la atribución. La más permisiva deja usar la obra incluso con fines comerciales, siempre que se dé crédito. Otra variante excluye el uso comercial. Otra obliga a que lo derivado circule con la misma licencia. Y otra impide además hacer obras derivadas. Elegir una sin leerla es el error más común.}
\milcpaso{3}{milcTurquesa}{Sobre IA y autoría hay tres consensos. Un modelo no es autor, porque la ley reconoce como autores a personas. Si la intervención humana es sustancial --- estructura, voz, edición, decisiones --- la obra es de quien la hizo. Y declarar el uso de la herramienta es práctica profesional aceptada, aunque en Colombia la ley todavía no lo exija.}
\milcpaso{4}{milcTurquesa}{Algoritmo para elegir licencia: (1) si tu objetivo es que el libro llegue a mucha gente, ve a una licencia abierta con atribución; (2) si quieres difundir pero no que otros lucren, añade la restricción comercial; (3) si quieres que lo derivado siga siendo abierto, usa la que obliga a mantener la licencia; (4) si piensas publicar comercialmente, reserva los derechos.}

\begin{drawbox}{milcTurquesa}{Anatomía de los derechos de una obra}
\textbf{Morales:} irrenunciables, intransferibles y perpetuos. \\[1mm] \textbf{Patrimoniales:} transferibles y negociables. \\[1mm] \textbf{Licencias abiertas:} todas parten de la atribución y varían en lo comercial y lo derivado. \\[1mm] \textbf{IA:} el modelo no es autor; la obra es de quien decidió. \\[1mm] \textbf{Declaración:} qué hizo la herramienta y qué hiciste tú.
\end{drawbox}

\begin{softbox}{milcMostaza}{milcGris}{Errores comunes y su corrección}
(1) No declarar el uso de IA, que en el trabajo editorial se considera una falta. (2) Confundir morales con patrimoniales y creer que se cede todo. (3) Elegir una licencia sin haberla leído. (4) Usar texto o imágenes ajenas sin atribuir la fuente. (5) Declarar el uso de IA con una frase genérica que no dice qué se hizo con ella.
\end{softbox}

\begin{infoband}{milcTurquesa}


\textbf{Lo que va al cuaderno (Actividad 2 · ANALIZA).} \textbf{Título:} «Actividad 2 --- Morales, patrimoniales y licencias». \textbf{Formato:} los dos tipos de derecho con ejemplo, las licencias con lo que permite cada una, los tres consensos sobre IA y la objeción que te puso tu pareja a la licencia elegida. \textbf{Extensión:} media página. \textbf{Sabes que terminaste cuando} podrías explicarle a alguien de tu casa qué licencia elegiste y por qué.
\end{infoband}

\puente{Tienes el método. Toca aplicarlo. Analizas tres casos recientes, escribes la ficha de derechos de tu libro y justificas la licencia que elegiste.}

\milcestacion{milcMagenta}{Estación 3 de 3 · Hacer}

\begin{softbox}{milcMagenta}{milcGris}{Cómo vas a construir tu producto}
\textbf{Actividad 3 · CREA --- La ficha de derechos de mi libro} (40 min · individual). (1) Busca tres casos recientes de disputa entre IA y derechos de autor. (2) Resume cada uno en un párrafo: quién demanda a quién, qué se discute y en qué va. (3) Escribe la ficha de derechos de tu libro con autoría, licencia, uso de IA, fuentes atribuidas y contacto. (4) En la declaración de IA di qué modelos usaste y en qué partes del proceso. (5) Justifica la licencia en cinco líneas. (6) Comprueba que la ficha se entienda sin que tú la expliques.
\end{softbox}

\milcpaso{1}{milcMagenta}{Los tres casos tienen que ser recientes, porque el campo cambia rápido y una sentencia de hace cinco años puede estar superada. Resúmelos con quién, qué se discute y en qué va. Si un caso sigue abierto, dilo: un pleito sin fallo no es una respuesta.}
\milcpaso{2}{milcMagenta}{La ficha de derechos tiene cinco piezas: autoría con nombre completo y año, licencia elegida y justificada, declaración del uso de IA, atribución de las fuentes que citaste y un contacto. Es la página de créditos de tu libro, escrita por ti.}
\milcpaso{3}{milcMagenta}{Sigue el patrón «transparencia antes que opacidad». Aunque la ley no lo exija, declarar el uso de la herramienta fortalece lo que firmas. Y la declaración tiene que ser específica: qué modelos, en qué partes, qué decidiste tú. «Se usó inteligencia artificial» no declara nada.}
\milcpaso{4}{milcMagenta}{Algoritmo: (1) busca los tres casos; (2) resume cada uno en un párrafo; (3) escribe la ficha con sus cinco piezas; (4) justifica la licencia en cinco líneas; (5) dásela a leer a alguien que no sepa de tu libro y comprueba si entiende qué puede y qué no puede hacer con él.}

\begin{drawbox}{milcMagenta}{Checklist antes de cerrar la ficha}
\checkbox Tres casos recientes resumidos, con su estado actual. \\[1mm] \checkbox Autoría con nombre completo y año. \\[1mm] \checkbox Licencia elegida y justificada en cinco líneas. \\[1mm] \checkbox Declaración de IA que dice qué modelos y en qué partes. \\[1mm] \checkbox Fuentes citadas atribuidas una por una. \\[1mm] \checkbox La ficha se entiende sin que tú la expliques.
\end{drawbox}

\begin{softbox}{milcMagenta}{milcGris}{Plantilla de trabajo}
\textbf{Casos.} \textit{Tres, con quién, qué se discute y en qué va.} \\[1mm] \textbf{Autoría.} \textit{Nombre completo y año.} \\[1mm] \textbf{Licencia.} \textit{Cuál y por qué, en cinco líneas.} \\[1mm] \textbf{Uso de IA.} \textit{Qué modelos, en qué partes, qué decidiste tú.} \\[1mm] \textbf{Fuentes.} \textit{Lo citado, atribuido.} \\[1mm] \textbf{Contacto.} \textit{Cómo escribirte.}
\end{softbox}

\begin{infoband}{milcMagenta}


\textbf{Lo que va al cuaderno (Actividad 3 · CREA).} \textbf{Título:} «Actividad 3 --- Los derechos de mi libro». \textbf{Formato:} los tres casos resumidos, la ficha completa y la justificación de la licencia. \textbf{Extensión:} una página. \textbf{Sabes que terminaste cuando} alguien que no conoce tu libro entiende, leyendo la ficha, qué puede y qué no puede hacer con él.
\end{infoband}

\puente{Lo que sigue resume \emph{qué} entregas hoy: los tres casos y la ficha de derechos. El criterio principal es que la declaración de uso de IA diga qué hiciste tú y qué hizo la herramienta, sin quedarse en una frase de compromiso.}

\milcestacion{milcMostaza}{Producto de la sesión}
\textbf{3 casos recientes analizados y ficha de derechos con licencia justificada}.
\begin{softbox}{milcMostaza}{milcGris}{Criterios de calidad}
(1) Tres casos recientes resumidos, con su estado actual declarado. (2) Autoría con nombre completo y año. (3) Licencia elegida y justificada, no copiada. (4) Declaración de IA que dice qué modelos y en qué partes del proceso. (5) Fuentes citadas atribuidas una por una. (6) La ficha se entiende sin explicación adicional.
\end{softbox}

\puente{Antes de cerrar, mira los derechos desde las cinco dimensiones humanas. Firmar con claridad es respeto por quien lee y por quienes escribieron antes que tú.}

\milcestacion{milcVino}{Evaluación liberadora · cinco dimensiones}

\begin{softbox}{milcVino}{milcGris}{Sentido de la evaluación}
Evaluar no es solo revisar una nota. Es reconocer qué aprendí, qué cambió en mí, qué emociones manejé, cómo me relaciono con la ciudad y la comunidad, y qué le dejaré a quien viene detrás.
\end{softbox}

\textbf{Mi reflexión liberadora en cinco dimensiones.} Completa cada frase iniciada:

\small
\begin{tabularx}{\textwidth}{>{\bfseries\RaggedRight\arraybackslash}p{3.4cm}Y>{\RaggedRight\arraybackslash}p{4.6cm}}
\rowcolor{milcVino}\color{white}Dimensión & \color{white}\bfseries Mi respuesta (frame starter) & \color{white}\bfseries Algo que descubrí\\
1. Personal & Lo que más cambió en mí fue \linead{6.5cm} & \linead{4.2cm}\\
2. Emocional & La emoción más fuerte fue \linead{2.5cm} y la regulé \linead{3cm} & \linead{4.2cm}\\
3. Ciudadana & En democracia esto importa porque \linead{6cm} & \linead{4.2cm}\\
4. Local (Cartago) & En mi barrio sería útil para \linead{6.5cm} & \linead{4.2cm}\\
5. Intergeneracional & A mi abuela/abuelo le diría \linead{6.5cm} & \linead{4.2cm}\\
\end{tabularx}
\normalsize

\begin{infoband}{milcVino}
\textbf{Para la próxima sustentación.} Vuelve a esta tabla en seis meses. Verás que las respuestas cambian — no porque lo aprendido cambie, sino porque tú cambias. La evaluación liberadora es una conversación contigo a través del tiempo.
\end{infoband}


\puente{Y para terminar, tres voces sobre de quién es lo que se hace. Dussel sobre quien queda fuera del sistema, Epicteto sobre lo que depende de uno, Floridi sobre no reducir a los demás a instrumentos.}

\milcestacion{milcGradeDark}{Triángulo de pensamiento · cierre filosófico}

\begin{softbox}{milcVino}{milcGris}{Dussel · lente del nosotros}
\textit{«El otro se revela realmente como otro… como el pobre, el oprimido; el que a la vera del camino, fuera del sistema, muestra su rostro sufriente y sin embargo desafiante: «¡Tengo hambre!, ¡tengo derecho a comer!».»}

{\footnotesize\itshape\color{milcNegro!70}--- Enrique Dussel · Filosofía de la liberación (1977), §2.4.4.2}

En el pleito por las piezas de oro discuten dos Estados. Los orfebres que las hicieron no aparecen: quedaron fuera del sistema que decide de quién son. Con los textos que entrenaron una herramienta pasa algo parecido, y la pregunta de quién queda fuera del reparto sigue abierta.

\textbf{De todo lo que hay detrás de mi libro, ¿quién no aparece nombrado en ninguna parte?}
\end{softbox}
\linead{\linewidth}\\[1mm]
\linead{\linewidth}

\begin{softbox}{milcMostaza}{milcGris}{Estoicismo · Epicteto · Enquiridión, 1 (c. 125 d.C.) · cuidado interior}
\textit{«Hay ciertas cosas que dependen de nosotros mismos, como la opinión, la inclinación, los deseos, la aversión, y en una palabra, todas nuestras operaciones. Otras hay también que no dependen, como el cuerpo, las riquezas, la reputación, los imperios.»}

{\footnotesize\itshape\color{milcNegro!70}--- Epicteto · Enquiridión, 1 (c. 125 d.C.)}

Los derechos morales son de lo que depende de ti y no se puede ceder aunque quieras. Los patrimoniales son de lo otro: se negocian, se venden, se pierden. Saber cuál es cuál evita firmar un contrato creyendo que entregas menos de lo que entregas.

\textbf{De lo que estoy declarando en mi ficha, ¿qué podría ceder mañana y qué seguirá siendo mío pase lo que pase?}
\end{softbox}
\linead{\linewidth}\\[1mm]
\linead{\linewidth}

\begin{softbox}{milcTurquesa}{milcGris}{Floridi · lente de la infoesfera}
\textit{«En una esfera pública plural, los demás no pueden reducirse a instrumentos, y hacen falta contención y respeto. (trad. propia, abreviada)»}

{\footnotesize\itshape\color{milcNegro!70}--- The Onlife Initiative (ed. Luciano Floridi) · The Onlife Manifesto (2015), § 4.3}

Citar es lo contrario de instrumentalizar. Cuando atribuyes una idea a quien la escribió, la tratas como obra de alguien y no como material disponible. Es la misma diferencia entre comprar una pieza y llevársela porque estaba ahí.

\textbf{¿Hay algo en mi libro que tomé de alguien y que todavía no está atribuido?}
\end{softbox}
\linead{\linewidth}\\[1mm]
\linead{\linewidth}

\begin{infoband}{milcNegro}
\textbf{Nota para el docente.} Las tres son citas textuales y llevan su referencia debajo. Puedes remitir a la obra en clase.
\end{infoband}

\milcestacion{milcVino}{Mi autoevaluación · marca una casilla por fila}

{\small
\setlength{\extrarowheight}{2.2pt}
\begin{tabularx}{\textwidth}{>{\RaggedRight\arraybackslash\bfseries}p{3.0cm}YYY}
\rowcolor{milcVino}\color{white}\bfseries Criterio & \color{white}\bfseries Lo logré & \color{white}\bfseries En proceso & \color{white}\bfseries Necesito apoyo\\[.6mm]
Comprensión del tema & \milcopcion{Explico las ideas con mis palabras.} & \milcopcion{Reconozco las ideas, falta precisión.} & \milcopcion{Necesito repasar lo básico.}\\[.8mm]
\rowcolor{milcGris}Pensamiento computacional & \milcopcion{Usé los pilares al armar mi producto.} & \milcopcion{Usé algunos pilares.} & \milcopcion{Necesito releer la estación 2.}\\[.8mm]
Aplicación y producto & \milcopcion{Mi producto cumple los criterios.} & \milcopcion{Está armado, con ajustes pendientes.} & \milcopcion{Aún está en construcción.}\\[.8mm]
\rowcolor{milcGris}Reflexión 5 dimensiones & \milcopcion{Respondí cada una con honestidad.} & \milcopcion{Respondí algunas con detalle.} & \milcopcion{Mis respuestas son muy generales.}\\[.8mm]
Triángulo de pensamiento & \milcopcion{Conecté las 3 voces con mi trabajo.} & \milcopcion{Conecté al menos una.} & \milcopcion{No las relaciono con mi proceso.}\\[.8mm]
\end{tabularx}}

\vspace{3mm}
\textbf{Hoy entendí que:}\\[1mm]
\linead{\linewidth}\\[1mm]
\linead{\linewidth}

\textbf{Mi compromiso para la próxima sesión es:}\\[1mm]
\linead{\linewidth}\\[1mm]
\linead{\linewidth}

\begin{softbox}{milcMostaza}{milcGris}{Cita de despedida · Marco Aurelio}
\textit{``El obstáculo en el camino se vuelve el camino. Lo que estorba al camino se vuelve camino.''}\\[1mm]
Cada error de hoy, bien observado, se convierte en pieza del camino que pisarás mañana.
\end{softbox}

\small
\begin{tabularx}{\textwidth}{>{\bfseries}p{3.6cm}Y}
\rowcolor{milcGradeDark!12}Nombre completo: & \linead{10cm}\\
Fecha: & \linead{4cm} \quad Curso/grupo: \linead{4cm}\\
Mi firma: & \linead{9cm}\\
\end{tabularx}
\normalsize

\begin{infoband}{milcGradeDark}
\textbf{Para la próxima sesión.} Llega con la ficha de derechos completa y los tres casos resumidos. En la sesión 7 vas a generar la portada del libro, y ahí la pregunta de los derechos vuelve, esta vez sobre imágenes.
\end{infoband}

\end{document}
